\clearpage
\section*{Problem 2}
\addcontentsline{toc}{section}{Problem 2}
\begingroup
\hypersetup{colorlinks=true,
            linkcolor=blue!50!black,
            citecolor=blue!50!black,
            urlcolor=blue!50!black}
\newcommand{\qtwoF}{F}
\newcommand{\qtwooo}{\mathfrak{o}}
\newcommand{\qtwopp}{\mathfrak{p}}
\newcommand{\qtwoqq}{\mathfrak{q}}
\newcommand{\qtwoC}{\mathbb{C}}
\newcommand{\qtwoZ}{\mathbb{Z}}
\newcommand{\qtwoGL}{\mathrm{GL}}
\newcommand{\qtwoMn}{M_n}
\newcommand{\qtwoSe}{\mathcal{S}_{\mathrm e}}
\newcommand{\qtwoSw}{\mathcal{S}}
\newcommand{\qtwoW}{\mathcal{W}}
\newcommand{\qtwovol}{\operatorname{vol}}
\newcommand{\qtwodiag}{\operatorname{diag}}
\newcommand{\qtwotr}{\operatorname{trace}}
\newcommand{\qtwottilde}[1]{\widetilde{#1}}
\newcommand{\qtwolRS}{\ell_{\mathrm{RS}}}
\newtheorem{qtwotheorem}{Theorem}
\newtheorem{qtwolemma}{Lemma}
\newtheorem{qtwoproposition}{Proposition}
\newtheorem{qtwocorollary}{Corollary}
\title{Solution to Problem 2: A Universal Test Vector\\
for Rankin--Selberg Integrals}
\author{}
\date{}
\maketitle


\section*{Statement and Answer}

Let $\qtwoF$ be a non-archimedean local field with ring of integers $\qtwooo$, maximal
ideal $\qtwopp$, residue cardinality $q$, and a nontrivial additive character
$\psi:\qtwoF\to\qtwoC^\times$ of conductor $\qtwooo$. Write $G_r:=\qtwoGL_r(\qtwoF)$, let $N_r<G_r$
be the upper-triangular unipotent subgroup, and embed $G_n\hookrightarrow
G_{n+1}$ as the upper-left block. Set $K_r:=\qtwoGL_r(\qtwooo)$, with Haar measures
giving $K_r$ and $N_r\cap K_r$ volume one.

\medskip
\noindent\textbf{Theorem.} \emph{Let $\Pi$ be a generic irreducible admissible
representation of $G_{n+1}$, realized in its $\psi^{-1}$-Whittaker model
$\qtwoW(\Pi,\psi^{-1})$. Then there exists a single $W\in\qtwoW(\Pi,\psi^{-1})$ with the
following property. For every generic irreducible admissible $\pi$ of $G_n$ with
conductor ideal $\qtwoqq$, generator $Q$ of $\qtwoqq^{-1}$, and
$u_Q:=I_{n+1}+QE_{n,n+1}$, there exists $V\in\qtwoW(\pi,\psi)$ such that}
\[
\int_{N_n\backslash G_n} W(\qtwodiag(g,1)u_Q)\,V(g)\,|\det g|^{s-\frac12}\,dg
\]
\emph{is finite and nonzero for all $s\in\qtwoC$.}

\medskip
\noindent The answer is \textbf{YES}. The single $W$ may be taken to be an
explicit, $\pi$-independent Kirillov-model vector, and for each $\pi$
the role of $V$ is played by the \emph{normalized newvector}. The integral is in
fact \emph{constant in $s$} and equal to an explicit nonzero scalar; this is why
finiteness and nonvanishing hold for all $s$, with no analytic continuation or
pole-cancellation required.

\medskip
\noindent\textbf{A failure mode to avoid.} One cannot choose $V$ as a
Schwartz bump in the Kirillov model making the integrand constant on its
support: $V$ has a central character, so it cannot be compactly supported modulo
the center. (Already for $n=1$ the relevant integral is a normalized Gauss sum,
whose integrand is genuinely non-constant.) Instead, the freedom is placed in
$W$: the chosen $W$ collapses the integral, after unfolding, to a single
$\det$-shell, whence it is constant in $s$; the value of that constant is a
\emph{Whittaker-model} $\varepsilon$-factor \textup{(}Gauss sum\textup{)}, which
is nonzero. We stress that the Godement--Jacquet functional equation, which
governs zeta integrals of \emph{matrix coefficients}, is \emph{not} available
here: a Whittaker function is not a matrix coefficient and $g\mapsto V(g^{-1})$
is not a Whittaker function, so the Godement--Jacquet dual cannot be formed from
$V$. The correct functional-equation input, when needed, is the
Jacquet--Piatetski-Shapiro--Shalika one on the Whittaker model.

\section*{Choice of $W$ and $V$}

For $x\in\qtwoF$ and $y\in\qtwoF^\times$ write
\[
d_y:=\qtwodiag(1,\dots,1,y)\in G_n\hookrightarrow G_{n+1},\qquad
u_x:=I_{n+1}+xE_{n,n+1}\in N_{n+1}.
\]

\noindent\textbf{The vector $W$.} Define $W_0$ on $G_n$ by
\begin{equation}\label{eq:W0}
W_0(g):=\int_{N_n}\mathbf 1_{K_n}(xg)\,\psi(x)\,dx,
\end{equation}
a function on $G_n$, left $(N_n,\psi^{-1})$-equivariant, which extends to an element of
$\qtwoW(\Pi,\psi^{-1})$ on $G_{n+1}$ by Lemma~\ref{lem:extend} below. We take
$W:=W_0$ to be this extension; it depends only on $\Pi,\psi,n$, \emph{not on $\pi$}, and the
problem's $\pi$-dependent shift $u_Q$ is applied as a right translation when forming
$g\mapsto W_0(\qtwodiag(g,1)u_Q)$. Throughout we identify $G_n$
with its image in $G_{n+1}$ under $g\mapsto\qtwodiag(g,1)$, and we use the same symbol
$W_0$ for the function on $G_n$ in \eqref{eq:W0} and for its chosen extension to
$G_{n+1}$; this is justified by Lemma~\ref{lem:extend}, which guarantees that the
extension restricts to \eqref{eq:W0} on $\qtwodiag(G_n,1)$.

The function $W_0$ on $G_n$ is supported on $N_nK_n$ and satisfies $W_0(k)=1$ for
$k\in K_n$. Indeed $W_0(g)\neq0$ forces $xg\in K_n$ for some $x\in N_n$, i.e.\
$g\in N_nK_n$; for $g=nk$ with $n\in N_n$, $k\in K_n$ the left $\psi^{-1}$-equivariance
gives $W_0(nk)=\psi^{-1}(n)W_0(k)$, and $W_0(k)=\int_{N_n}\mathbf 1_{K_n}(xk)\psi(x)\,dx
=\int_{N_n\cap K_n}\psi(x)\,dx=1$ because $\psi$ is trivial on $N_n\cap K_n$ (conductor
$\qtwooo$) and $\qtwovol(N_n\cap K_n)=1$.

\medskip
We now justify the two facts about $W_0$ used below: that \eqref{eq:W0} is the
restriction to $\qtwodiag(G_n,1)$ of a genuine element of $\qtwoW(\Pi,\psi^{-1})$, and that
this extension satisfies the right-$u_1$ translation identity. Let
$P_{n+1}:=\{p\in G_{n+1}:\text{last row of }p\text{ is }(0,\dots,0,1)\}$ be the
mirabolic subgroup of $G_{n+1}$; it is generated by $\qtwodiag(G_n,1)$ and $N_{n+1}$, and
$\qtwodiag(G_n,1)$ normalizes the abelian unipotent radical
$U:=\{I_{n+1}+v:\ v=\sum_{i=1}^{n}v_iE_{i,n+1}\}\cong\qtwoF^n$ of $P_{n+1}$.

\begin{qtwolemma}[Realizability in $\qtwoW(\Pi,\psi^{-1})$ via the Kirillov model]\label{lem:extend}
Let $\qtwoSw(N_n\backslash G_n;\psi^{-1})$ denote the space of locally constant functions
$f:G_n\to\qtwoC$ with $f(ng)=\psi^{-1}(n)f(g)$ for $n\in N_n$, $g\in G_n$, and with support
compact modulo $N_n$. Then\textup{:}
\begin{enumerate}[label=\textup{(\alph*)}]\setlength\itemsep{2pt}
\item \emph{\textup{(}Injectivity of mirabolic restriction.\textup{)}} The restriction map
$\rho:\qtwoW(\Pi,\psi^{-1})\to\{\text{functions on }P_{n+1}\}$, $W\mapsto W|_{P_{n+1}}$, is
injective\textup{;} equivalently $W\in\qtwoW(\Pi,\psi^{-1})$ is determined by its values on
$\qtwodiag(G_n,1)$.
\item \emph{\textup{(}Surjectivity onto compactly supported data.\textup{)}} For every
$f\in\qtwoSw(N_n\backslash G_n;\psi^{-1})$ there exists $W\in\qtwoW(\Pi,\psi^{-1})$ with
$W(\qtwodiag(g,1))=f(g)$ for all $g\in G_n$.
\end{enumerate}
In particular, the function $W_0$ of \eqref{eq:W0} lies in
$\qtwoSw(N_n\backslash G_n;\psi^{-1})$, so by \textup{(b)} there is a genuine
$W\in\qtwoW(\Pi,\psi^{-1})$ on all of $G_{n+1}$ with $W(\qtwodiag(g,1))=W_0(g)$ for all
$g\in G_n$\textup{;} by \textup{(a)} this $W$ is unique, and we denote it again by $W_0$.
\end{qtwolemma}

\begin{proof}
Throughout, $P_{n+1}$ is the mirabolic subgroup and
$U=\{I_{n+1}+\sum_{i=1}^n v_iE_{i,n+1}\}\cong\qtwoF^n$ its unipotent radical, so that
$P_{n+1}=\qtwodiag(G_n,1)\ltimes U$ as a semidirect product. The character $\psi^{-1}$ of
$N_{n+1}$ restricts to a character of $U$, namely $I_{n+1}+\sum_i v_iE_{i,n+1}\mapsto
\psi^{-1}(v_n)$, because among the entries $v_i$ \textup{(}all in column $n+1$\textup{)} only
$v_n$ sits on the superdiagonal.

\smallskip
\emph{Proof of \textup{(a)} \textup{(}injectivity\textup{)}.}
Every $W\in\qtwoW(\Pi,\psi^{-1})$ is left $(N_{n+1},\psi^{-1})$-equivariant, and the Iwasawa
decomposition $G_{n+1}=N_{n+1}A_{n+1}K_{n+1}$ \textup{(}with $A_{n+1}$ the diagonal torus\textup{)}
shows that $W$ is determined by its restriction to $A_{n+1}K_{n+1}$, indeed by its restriction to
any set meeting every $N_{n+1}$-orbit. The substance of \textup{(a)} is rather the standard
Bernstein--Zelevinsky fact that the restriction of the \emph{abstract} Whittaker space to
$P_{n+1}$ is already faithful as a $P_{n+1}$-module; we use it only in the weak form needed here:
distinct elements of $\qtwoW(\Pi,\psi^{-1})$ restrict to distinct functions on $P_{n+1}$. This is
\cite[2.5,~3.5]{Bernstein} \textup{(}see also \cite[(2.1)]{JPSS83}\textup{)}: the
$P_{n+1}$-module $\Pi|_{P_{n+1}}$ is the Kirillov model, on which $P_{n+1}$ acts faithfully, so the
$P_{n+1}$-equivariant map $\rho$ from the \textup{(}irreducible, hence faithful as a
$G_{n+1}$-module\textup{)} space $\qtwoW(\Pi,\psi^{-1})$ is injective. We record \textup{(a)} for
completeness; it is \textup{(b)} that is used in the sequel, while \textup{(a)} guarantees that the
``$W_0$'' produced by \textup{(b)} is unambiguous.

\smallskip
\emph{Proof of \textup{(b)} \textup{(}realizability of compactly supported data\textup{)}.}
Restricting Whittaker functions to $P_{n+1}$ and then to $\qtwodiag(G_n,1)$ realizes
$\Pi|_{P_{n+1}}$ as a space of functions on $N_n\backslash G_n$ that transform by $\psi^{-1}$ under
left translation by $N_n$; this is the \emph{Kirillov model} $\mathcal K(\Pi,\psi^{-1})$. The
identification is legitimate because $P_{n+1}=\qtwodiag(G_n,1)\ltimes U$ and $W$ transforms by the
character $\psi^{-1}|_U$ under right \textup{(}equivalently left, after the semidirect product
bookkeeping\textup{)} translation by $U$, so the values of $W$ on $P_{n+1}$ are recovered from its
values on $\qtwodiag(G_n,1)$ together with the fixed $U$-character; concretely, for $g\in G_n$ and
$u'=I_{n+1}+\sum_i v_iE_{i,n+1}\in U$ one has
$W(\qtwodiag(g,1)\,u')=\psi^{-1}\big((\qtwodiag(g,1)u'\qtwodiag(g,1)^{-1})_{n,n+1}\big)\,
W(\qtwodiag(g,1))$ by the same left-equivariance computation as in
Lemma~\ref{lem:u1translate} below. The Bernstein--Zelevinsky theorem on the Kirillov model now
states:
\begin{itemize}\setlength\itemsep{2pt}
\item[(K1)] \emph{\textup{(}Kirillov model contains the compactly supported functions;
\cite[\S\S2.5,~3.5]{Bernstein}, \cite[(2.3) and \S2]{JPSS83}.\textup{)}} The Kirillov space
$\mathcal K(\Pi,\psi^{-1})$ contains $\qtwoSw(N_n\backslash G_n;\psi^{-1})$. Equivalently, every
$f\in\qtwoSw(N_n\backslash G_n;\psi^{-1})$ is the restriction to $\qtwodiag(G_n,1)$ of some
$W\in\qtwoW(\Pi,\psi^{-1})$.
\end{itemize}
\textup{(}The mechanism behind (K1): the top Bernstein--Zelevinsky derivative $\Pi^{(n+1)}$ is one
dimensional for generic $\Pi$, and the Bernstein--Zelevinsky filtration of $\Pi|_{P_{n+1}}$ has a
bottom layer isomorphic to the compactly induced representation
$\mathrm{ind}_{N_{n+1}}^{P_{n+1}}\psi^{-1}$, whose space of functions on $N_n\backslash G_n$ is
exactly $\qtwoSw(N_n\backslash G_n;\psi^{-1})$. Genericity of $\Pi$ is what guarantees this bottom
layer is present, i.e.\ that the Whittaker/Kirillov model exists and contains all compactly
supported data.\textup{)} Fact (K1) is precisely assertion \textup{(b)}.

\smallskip
\emph{Verification that $W_0\in\qtwoSw(N_n\backslash G_n;\psi^{-1})$.}
$\bullet$ \emph{Left equivariance:} for $n_0\in N_n$ the substitution $x\mapsto xn_0^{-1}$ in
\eqref{eq:W0} \textup{(}note $N_n$ is unimodular and $\psi$ is a character of $N_n$\textup{)} gives
$W_0(n_0 g)=\int_{N_n}\mathbf 1_{K_n}(xn_0g)\psi(x)\,dx=\int_{N_n}\mathbf 1_{K_n}(xg)\psi(xn_0^{-1})\,dx
=\psi(n_0^{-1})W_0(g)=\psi^{-1}(n_0)W_0(g)$.
$\bullet$ \emph{Smoothness:} for $k'\in K_n$ we have $\mathbf 1_{K_n}(xgk')=\mathbf 1_{K_n}(xg)$, hence
$W_0(gk')=W_0(g)$; thus $W_0$ is right $K_n$-invariant, in particular locally constant.
$\bullet$ \emph{Compact support mod $N_n$:} $W_0(g)\neq0$ forces $xg\in K_n$ for some $x\in N_n$,
i.e.\ $g\in N_nK_n$; the image of $N_nK_n$ in $N_n\backslash G_n$ is
$(N_n\cap K_n)\backslash K_n$, which is compact. Hence $W_0\in\qtwoSw(N_n\backslash G_n;\psi^{-1})$,
and (K1) applies to give the required $W$.
\end{proof}

\begin{qtwolemma}[Off-mirabolic right-translation identity]\label{lem:u1translate}
Let $W\in\qtwoW(\Pi,\psi^{-1})$ be arbitrary, $u_1:=I_{n+1}+E_{n,n+1}\in N_{n+1}$, and
$g\in G_n$ identified with $\qtwodiag(g,1)\in P_{n+1}$. Then for every $h\in G_{n+1}$,
\[
W(\qtwodiag(g,1)\,u_1\,h)=\psi(-g_{nn})\,W(\qtwodiag(g,1)\,h).
\]
In particular, applying this with $h=d_Q^{-1}$ to the extension $W=W_0$ of
Lemma~\ref{lem:extend}, and using $t_Q=u_1 d_Q^{-1}$,
\[
W_0(g\,t_Q)=W_0(\qtwodiag(g,1)\,u_1\,d_Q^{-1})=\psi(-g_{nn})\,W_0(\qtwodiag(g,1)\,d_Q^{-1})
=\psi(-g_{nn})\,W_0(g\,d_Q^{-1}),
\]
the last equality because $\qtwodiag(g,1)\,d_Q^{-1}=\qtwodiag(g\,d_Q^{-1},1)$ lies in
$\qtwodiag(G_n,1)$, where the extension restricts to \eqref{eq:W0}.
\end{qtwolemma}

\begin{proof}
The mechanism is that right multiplication by the off-diagonal unipotent $u_1$ is converted, by
conjugation, into \emph{left} multiplication by an element of $N_{n+1}$, on which $W$ transforms by
the generic character $\psi^{-1}$. Set
\[
n_g:=\qtwodiag(g,1)\,u_1\,\qtwodiag(g,1)^{-1}\in G_{n+1},
\qquad\text{so that}\qquad \qtwodiag(g,1)\,u_1=n_g\,\qtwodiag(g,1).
\]

\emph{Step 1: $n_g\in N_{n+1}$ and its entries.}
Write $e_1,\dots,e_{n+1}$ for the standard column basis and $M:=\qtwodiag(g,1)$. Since
$E_{n,n+1}=e_n\,e_{n+1}^{\mathsf T}$, we have
\[
n_g=M\,(I_{n+1}+E_{n,n+1})\,M^{-1}=I_{n+1}+M\,E_{n,n+1}\,M^{-1}
=I_{n+1}+(M e_n)\,(e_{n+1}^{\mathsf T}M^{-1}).
\]
Now $M e_n=\sum_{i=1}^{n}g_{in}e_i$ \textup{(}the $n$th column of $\qtwodiag(g,1)$, which is the
$n$th column of $g$ in the first $n$ coordinates and $0$ in the last\textup{)}, while
$e_{n+1}^{\mathsf T}M^{-1}=e_{n+1}^{\mathsf T}$ because $M$ fixes the last row, hence so does $M^{-1}$.
Therefore
\[
n_g=I_{n+1}+\Big(\sum_{i=1}^n g_{in}e_i\Big)e_{n+1}^{\mathsf T}
=I_{n+1}+\sum_{i=1}^{n}g_{in}\,E_{i,n+1}.
\]
This is an upper-triangular unipotent matrix \textup{(}all added entries lie in column $n+1$,
strictly above the diagonal\textup{)}, so $n_g\in N_{n+1}$, with nonzero off-diagonal entries
exactly $(n_g)_{i,n+1}=g_{in}$ for $1\le i\le n$.

\emph{Step 2: evaluating the generic character.}
The generic character of $N_{n+1}$ attached to $\psi$ is
$u\mapsto\psi\big(\sum_{i=1}^{n}u_{i,i+1}\big)$, depending only on the superdiagonal entries
$u_{i,i+1}$. Among the nonzero entries $(n_g)_{i,n+1}=g_{in}$ \textup{(}$1\le i\le n$\textup{)},
the only index of the form $(i,i+1)$ is $i=n$, giving the superdiagonal entry $(n_g)_{n,n+1}=g_{nn}$;
the entries $(i,n+1)$ with $i<n$ satisfy $n+1>i+1$, hence lie \emph{strictly above} the
superdiagonal and do not enter the generic character. Therefore the character $\psi^{-1}$ of
$N_{n+1}$ \textup{(}the one for which $\qtwoW(\Pi,\psi^{-1})$ is the $\psi^{-1}$-Whittaker
model\textup{)} takes the value
\[
\psi^{-1}(n_g)=\psi\!\big(-(n_g)_{n,n+1}\big)=\psi(-g_{nn}).
\]

\emph{Step 3: applying left equivariance.}
For any $h\in G_{n+1}$, using $\qtwodiag(g,1)u_1=n_g\,\qtwodiag(g,1)$ and the left
$(N_{n+1},\psi^{-1})$-equivariance $W(u\,x)=\psi^{-1}(u)W(x)$ \textup{(}$u\in N_{n+1}$\textup{)},
\[
W(\qtwodiag(g,1)\,u_1\,h)=W(n_g\,\qtwodiag(g,1)\,h)
=\psi^{-1}(n_g)\,W(\qtwodiag(g,1)\,h)=\psi(-g_{nn})\,W(\qtwodiag(g,1)\,h).
\]
This is the asserted identity.

\emph{The displayed consequence.}
We verify $t_Q=u_1 d_Q^{-1}=d_Q^{-1}u_Q$. Here $d_Q^{-1}=\qtwodiag(1,\dots,1,Q^{-1},1)$ has
$Q^{-1}$ in the $(n,n)$ slot, and $u_Q=I_{n+1}+QE_{n,n+1}$. Both products are unipotent with their
unique nonzero off-diagonal entry at position $(n,n+1)$: in $u_1 d_Q^{-1}$ that entry is
$(u_1)_{n,n+1}\cdot(d_Q^{-1})_{n+1,n+1}=1\cdot1=1$, and in $d_Q^{-1}u_Q$ it is
$(d_Q^{-1})_{nn}\cdot Q=Q^{-1}\cdot Q=1$; all diagonal entries agree, so the two products are equal,
call it $t_Q$ \textup{(}with $(t_Q)_{n,n+1}=1$\textup{)}. Taking $h=d_Q^{-1}$ in the identity and
$W=W_0$ gives
\[
W_0(\qtwodiag(g,1)\,t_Q)=W_0(\qtwodiag(g,1)\,u_1\,d_Q^{-1})=\psi(-g_{nn})\,W_0(\qtwodiag(g,1)\,d_Q^{-1}),
\]
and $\qtwodiag(g,1)\,d_Q^{-1}=\qtwodiag(g d_Q^{-1},1)\in\qtwodiag(G_n,1)$, where $W_0$ restricts to
\eqref{eq:W0}; this yields the displayed formula
$W_0(g\,t_Q)=\psi(-g_{nn})\,W_0(g\,d_Q^{-1})$.
\end{proof}

\noindent\textbf{The vector $V$.} Given $\pi$ with conductor $\qtwoqq=\qtwopp^{c}$
\textup{(}$c\ge0$\textup{)}, fix a generator $Q$ of $\qtwoqq^{-1}$, so that
$|Q|=[\qtwooo:\qtwoqq]=q^{c}\ge1$. We will choose $V\in\qtwoW(\pi,\psi)$ depending on $\pi$
\textup{(}as the problem permits\textup{)}; the load-bearing point is that $W=W_0$ does
\emph{not} depend on $\pi$. For the explicit evaluation we take $V$ to be the normalized
newvector of $\pi$: the unique $K_1(\qtwoqq)$-invariant vector with $V(1)=1$, where
$K_1(\qtwoqq)=\{g\in K_n:\text{last row}\equiv(0,\dots,0,1)\bmod\qtwoqq\}$
\cite{JPSS81,Matringe}. Recall the conductor is characterized by the $\varepsilon$-factor:
\begin{equation}\label{eq:eps}
\varepsilon(\tfrac12+s,\pi,\psi)=|Q|^{-s}\,\varepsilon(\tfrac12,\pi,\psi),
\qquad |\varepsilon(\tfrac12,\pi,\psi)|=1 .
\end{equation}

\medskip
\noindent\textbf{Uniformity: the role of the single, $\pi$-independent $W_0$.}
We emphasize the logical order, which is the crux of the problem. The vector
$W_0\in\qtwoW(\Pi,\psi^{-1})$ is constructed once and for all in \eqref{eq:W0} and
Lemma~\ref{lem:extend}, \emph{before any $\pi$ is named}: its definition refers only to
$\Pi$, $\psi$ and $n$, never to $\pi$, to the conductor $c$, or to $Q$. When a representation
$\pi$ of $G_n$ is subsequently given, the only $\pi$-dependent object entering the
Rankin--Selberg integral is the group element
$u_Q=I_{n+1}+QE_{n,n+1}\in N_{n+1}$, which acts on the \emph{fixed} $W_0$ by right
translation, $W_0\mapsto u_QW_0=:\big(g\mapsto W_0(\qtwodiag(g,1)u_Q)\big)$. Thus a single
$W_0$ serves \emph{all} $\pi$ simultaneously, and there is no boundedness hypothesis on the
conductor anywhere: the analysis below holds verbatim for every $c\ge0$, hence for arbitrarily
ramified $\pi$.

\section*{Uniformity: the structure of $g\mapsto W_0(\qtwodiag(g,1)u_Q)$ is independent of $a(\pi)$}

The heart of the present issue is to show \emph{rigorously} that the single, $\pi$-independent
vector $W_0$ of Lemma~\ref{lem:extend} retains, after the $\pi$-dependent right translation by
$u_Q$, exactly the support and equivariance properties required by the test-vector argument, and
that it does so \emph{uniformly in the conductor exponent} $c=a(\pi)$, which is unbounded. We
isolate this as a self-contained lemma; everything afterwards is a formal consequence.

\begin{qtwolemma}[Uniform structure of the $u_Q$-translate]\label{lem:uniform}
Let $W_0\in\qtwoW(\Pi,\psi^{-1})$ be the fixed vector of Lemma~\ref{lem:extend}, defined
\emph{before any $\pi$ is named} and depending only on $\Pi,\psi,n$. For every generic
irreducible admissible $\pi$ of $G_n$, with conductor $\qtwoqq=\qtwopp^{c}$ \textup{(}any
$c\ge0$\textup{)}, generator $Q$ of $\qtwoqq^{-1}$ \textup{(}so $|Q|=q^{c}$\textup{)}, and
$u_Q=I_{n+1}+QE_{n,n+1}$, define $\Phi_Q:G_n\to\qtwoC$ by
\[
\Phi_Q(g):=W_0(\qtwodiag(g,1)\,u_Q).
\]
Then, for \emph{every} $c\ge0$, the function $\Phi_Q$ satisfies the following three properties,
in which the conductor enters \emph{only} through the modulus-one phase $\psi(-Q\,g_{nn})$ and
through nothing else\textup{:}
\begin{enumerate}[label=\textup{(\roman*)}]\setlength\itemsep{2pt}
\item \emph{\textup{(}Closed form\textup{)}} $\Phi_Q(g)=\psi(-Q\,g_{nn})\,W_0(g)$ for all
$g\in G_n$, where $W_0$ on $\qtwodiag(G_n,1)$ is given by \eqref{eq:W0}.
\item \emph{\textup{(}$c$-independent support\textup{)}} $\operatorname{supp}\Phi_Q\subseteq N_nK_n$,
the \emph{same} set for every $c$; in particular $\Phi_Q$ is supported on a single $\det$-shell,
$|\det g|=1$ on $\operatorname{supp}\Phi_Q$.
\item \emph{\textup{(}Twisted left equivariance\textup{)}} for $n_0\in N_n$ and $g\in G_n$,
\[
\Phi_Q(n_0 g)=\psi^{-1}(n_0)\,\Phi_Q(g),
\]
the \emph{same} equivariance as $W_0$ itself, for every $c$.
\end{enumerate}
Consequently the pairing $g\mapsto\Phi_Q(g)V(g)|\det g|^{s-1/2}$ that defines the
Rankin--Selberg integral is, for every $\pi$ and every $c\ge0$, an $N_n$-invariant\footnote{The
twist $\psi^{-1}$ of $\Phi_Q$ cancels the twist $\psi$ of $V$, so the product descends to
$N_n\backslash G_n$ as required for the integral to make sense.} function on
$N_n\backslash G_n$ supported on the fixed compact set $(N_n\cap K_n)\backslash K_n$, with no
dependence of the support on $c$.
\end{qtwolemma}

\begin{proof}
\emph{(i).} Write $M:=\qtwodiag(g,1)$. As computed in Lemma~\ref{lem:u1translate}
\textup{(}Step~1, with the off-diagonal entry scaled by $Q$\textup{)},
\[
n_g^{(Q)}:=M\,u_Q\,M^{-1}=I_{n+1}+Q\sum_{i=1}^{n}g_{in}\,E_{i,n+1}\in N_{n+1},
\]
so that $M\,u_Q=n_g^{(Q)}\,M$. The only superdiagonal entry of $n_g^{(Q)}$ is
$(n_g^{(Q)})_{n,n+1}=Q\,g_{nn}$ \textup{(}the entries $(i,n+1)$ with $i<n$ lie strictly above the
superdiagonal and do not enter the generic character\textup{)}, so
$\psi^{-1}(n_g^{(Q)})=\psi(-Q\,g_{nn})$. By the left $(N_{n+1},\psi^{-1})$-equivariance of $W_0$,
\[
\Phi_Q(g)=W_0(M\,u_Q)=W_0(n_g^{(Q)}\,M)=\psi^{-1}(n_g^{(Q)})\,W_0(M)=\psi(-Q\,g_{nn})\,W_0(g),
\]
where $W_0(M)=W_0(\qtwodiag(g,1))$ equals \eqref{eq:W0} by Lemma~\ref{lem:extend}. This is the
closed form, valid for every $c$; the conductor $c$ enters only through the unimodular scalar
$\psi(-Q\,g_{nn})$.

\emph{(ii).} By (i), $\Phi_Q(g)\ne0$ implies $W_0(g)\ne0$ \textup{(}the phase is never zero\textup{)},
and conversely $|\Phi_Q(g)|=|W_0(g)|$. Hence
$\operatorname{supp}\Phi_Q=\operatorname{supp}W_0\subseteq N_nK_n$ by the support computation in
Lemma~\ref{lem:extend}; this set is independent of $\pi,Q,c$. On $N_nK_n$ one has
$|\det g|=|\det k|=1$ for $g=n k$ \textup{(}$n\in N_n$, $k\in K_n$\textup{)}, so $\Phi_Q$ is
supported on the single shell $|\det g|=1$.

\emph{(iii).} For $n_0\in N_n\subset N_{n+1}$, the $(n,n)$-entry is unchanged under left
multiplication by $n_0$ \textup{(}elements of $N_n$ are upper unipotent, so $(n_0 g)_{nn}=g_{nn}$
because the last row of $n_0$ is $e_n$\textup{)}, whence the phase $\psi(-Q\,g_{nn})$ is
left-$N_n$-invariant. Combined with the left $\psi^{-1}$-equivariance of $W_0$
\textup{(}Lemma~\ref{lem:extend}\textup{)},
\[
\Phi_Q(n_0 g)=\psi(-Q\,(n_0g)_{nn})\,W_0(n_0 g)
=\psi(-Q\,g_{nn})\,\psi^{-1}(n_0)\,W_0(g)=\psi^{-1}(n_0)\,\Phi_Q(g).
\]
This holds for every $c$.

\emph{Final assertion.} The function $g\mapsto\Phi_Q(g)V(g)$ has left $N_n$-equivariance
$\psi^{-1}(n_0)\cdot\psi(n_0)=1$ by (iii) and the $\psi$-equivariance of $V$, so it is
$N_n$-invariant and descends to $N_n\backslash G_n$; and $|\det g|^{s-1/2}$ is $N_n$-invariant.
By (ii) the integrand is supported on the image of $K_n$, the compact set
$(N_n\cap K_n)\backslash K_n$, for every $c$. This is precisely the ``support and equivariance,
uniformly in $a(\pi)$'' required.
\end{proof}

\noindent\emph{Reading of the issue.} Lemma~\ref{lem:uniform} settles the conceptual core flagged
in the issue: the conductor dependence is carried \emph{entirely} by the scalar phase
$\psi(-Q\,g_{nn})$ \textup{(}of modulus $1$\textup{)}, while the support
$\operatorname{supp}\Phi_Q=N_nK_n$ and the $\psi^{-1}$-equivariance of $\Phi_Q$ are
\emph{literally the same} for every $\pi$, regardless of how large $a(\pi)=c$ is. There is no
hidden $\pi$- or $c$-dependence in the geometry of the translated vector; the only $\pi$-dependent
ingredient is a bounded twist. This is exactly why one fixed $W_0$, chosen before any $\pi$,
serves all $\pi$ simultaneously and across all conductors.

\section*{Collapse to a single shell, and $s$-independence}

We now compute the problem's integral with the fixed vector $W=W_0$ and verify that it
collapses, \emph{uniformly in the conductor $c$}, to an $s$-independent constant. This step is
entirely elementary \textup{(}one application of the off-mirabolic translation identity, then an
unfolding\textup{)} and uses \emph{no} functional equation.

\begin{qtwolemma}\label{lem:unfold}
Let $W_0\in\qtwoW(\Pi,\psi^{-1})$ be the fixed vector of Lemma~\ref{lem:extend}. For every
generic irreducible admissible $\pi$ of $G_n$, with conductor $\qtwoqq=\qtwopp^{c}$, generator
$Q$ of $\qtwoqq^{-1}$, $u_Q=I_{n+1}+QE_{n,n+1}$, every $V\in\qtwoW(\pi,\psi)$, and every
$s\in\qtwoC$,
\[
\int_{N_n\backslash G_n} W_0(\qtwodiag(g,1)\,u_Q)\,V(g)\,|\det g|^{s-\frac12}\,dg
=\int_{G_n}\psi(-Q\,g_{nn})\,\mathbf 1_{K_n}(g)\,V(g)\,|\det g|^{s-\frac12}\,dg ,
\]
and the right-hand side is supported on the single shell $|\det g|=1$. In particular the
integral converges absolutely \textup{(}the domain is the compact set $K_n$\textup{)} and equals
the $s$-independent constant
\begin{equation}\label{eq:Idef}
I:=\int_{K_n}\psi(-Q\,k_{nn})\,V(k)\,dk\qquad\textup{(Haar on $K_n$ of total mass $1$),}
\end{equation}
for \emph{every} conductor $c\ge0$.
\end{qtwolemma}

\begin{proof}
\emph{Step 1 \textup{(}closed form of the $u_Q$-translate, uniform in $c$\textup{)}.}
By Lemma~\ref{lem:uniform}(i),
\[
W_0(\qtwodiag(g,1)\,u_Q)=\psi(-Q\,g_{nn})\,W_0(g),
\]
where $W_0(g)$ denotes the value on $\qtwodiag(g,1)\in\qtwodiag(G_n,1)$, equal to \eqref{eq:W0}
by Lemma~\ref{lem:extend}. Recall the mechanism \textup{(}proof of Lemma~\ref{lem:uniform}\textup{)}:
$u_Q$ lies in the abelian unipotent radical $U$ of $P_{n+1}$, conjugation by $M:=\qtwodiag(g,1)$
turns the right multiplication by $u_Q$ into left multiplication by
$n_g^{(Q)}=I_{n+1}+Q\sum_{i=1}^{n}g_{in}E_{i,n+1}\in N_{n+1}$, whose unique superdiagonal entry is
$Q\,g_{nn}$, and the left $(N_{n+1},\psi^{-1})$-equivariance of $W_0$ produces the scalar
$\psi^{-1}(n_g^{(Q)})=\psi(-Q\,g_{nn})$. This identity holds verbatim for every conductor
exponent $c\ge0$; $c$ enters only through the modulus-one phase $\psi(-Q\,g_{nn})$.

\emph{Step 2 \textup{(}unfolding\textup{)}.}
By \eqref{eq:W0}, $W_0(g)=\int_{N_n}\mathbf 1_{K_n}(x\,g)\,\psi(x)\,dx$. Substituting Step 1 into
the integral and unfolding $N_n\backslash G_n$ against $N_n$ \textup{(}the map
$(x,N_ng)\mapsto xg$ identifies $N_n\times(N_n\backslash G_n)$ with $G_n$ as measure spaces; we
write $g=x^{-1}g'$, $g'=xg$\textup{)} we use three invariances. First, the $(n,n)$-entry is
unchanged by left multiplication by $N_n$ \textup{(}whose elements have last row $e_n$\textup{)},
so $(x^{-1}g')_{nn}=g'_{nn}$ and the phase $\psi(-Q\,g_{nn})$ becomes $\psi(-Q\,g'_{nn})$.
Second, the left $\psi$-equivariance of $V$ gives
$V(x^{-1}g')=\psi(x^{-1})V(g')=\psi(-x)V(g')$, where $\psi(x)$ here denotes the value of the
superdiagonal character of $N_n$ at $x$; this $\psi(-x)$ cancels the phase $\psi(x)$ coming from
\eqref{eq:W0} \textup{(}the cancellation $\psi(x)\psi(x^{-1})=1$ holds for every $n$ and is
independent of $c$\textup{)}. Third, $|\det(x^{-1}g')|=|\det g'|$. The result is
\[
\int_{G_n}\psi(-Q\,g'_{nn})\,\mathbf 1_{K_n}(g')\,V(g')\,|\det g'|^{s-\frac12}\,dg' ,
\]
which is the asserted $G_n$-integral.

\emph{Step 3 \textup{(}single-shell collapse\textup{)}.}
The integrand is supported where $g'\in K_n$, on which $|\det g'|=1$, so $|\det g'|^{s-\frac12}=1$
\emph{for all $s$}. Hence the integral equals
$\int_{K_n}\psi(-Q\,k_{nn})\,V(k)\,dk=I$, independent of $s$, and the domain $K_n$ is compact, so
the integral converges absolutely. None of Steps~1--3 involves any restriction on $c$: the
conductor enters only through the bounded phase $\psi(-Q\,k_{nn})$ \textup{(}of modulus $1$\textup{)},
so the conclusion is uniform in $c\ge0$.
\end{proof}

\begin{qtwoproposition}[The single-shell collapse, and why no multi-shell function arises]\label{prop:singleshell}
With the notation of Lemma~\ref{lem:unfold}, the integrand
$g\mapsto W_0(\qtwodiag(g,1)u_Q)\,V(g)\,|\det g|^{s-1/2}$, viewed as a function on
$N_n\backslash G_n$, is supported on the \emph{single} determinant shell
$S_1:=\{g\in N_n\backslash G_n:|\det g|=1\}$, for every $\pi$ and every conductor $c\ge0$.
Consequently, writing the integral as a sum over shells $S_{q^{m}}=\{|\det g|=q^{m}\}$,
\[
\qtwolRS(s,W_0,V)=\sum_{m\in\qtwoZ}\Big(\int_{S_{q^m}}W_0(\qtwodiag(g,1)u_Q)\,V(g)\,dg\Big)\,q^{-m(s-1/2)},
\]
\emph{only the term $m=0$ is nonzero}, so $\qtwolRS(s,W_0,V)=I\cdot q^{0}=I$ is a single monomial
in $q^{\pm s}$ \textup{(}indeed the constant monomial $c^{\,s-1/2}$ with $c=1$\textup{)}, hence
entire in $s$ and nowhere vanishing as a function of $s$ precisely when $I\ne0$.
\end{qtwoproposition}

\begin{proof}
By Lemma~\ref{lem:uniform}(ii) the function $g\mapsto W_0(\qtwodiag(g,1)u_Q)$ has support contained
in $N_nK_n$, and on $N_nK_n$ one has $|\det g|=1$. Hence $W_0(\qtwodiag(g,1)u_Q)=0$ whenever
$|\det g|\ne1$, so the shell integral over $S_{q^m}$ vanishes for every $m\ne0$, regardless of
$V$. The displayed shell decomposition is the partition of $N_n\backslash G_n$ by the
locally constant function $|\det\,\cdot\,|$ \textup{(}which takes values in $q^{\qtwoZ}$\textup{)},
on each shell of which $|\det g|^{s-1/2}$ is the constant $q^{-m(s-1/2)}$; the integral converges
absolutely there because the integrand is compactly supported \textup{(}mod $N_n$\textup{)} on the
single shell $S_1$. The $m=0$ coefficient is $\int_{S_1}W_0(\qtwodiag(g,1)u_Q)V(g)\,dg=I$ by
Lemma~\ref{lem:unfold}.
\end{proof}

\noindent\emph{Why the documented multi-shell danger does not occur here.} The fatal failure mode
recorded in the issue is the following: if one takes $W$ \textup{(}or $V$\textup{)} to be a
\emph{genuine} Whittaker/newvector whose values on the diagonal torus are governed by the
recursive Casselman--Shalika / Bernstein--Zelevinsky formula, then those values are supported on
an \emph{unbounded range} of torus cosets, and the Rankin--Selberg integral reproduces
$L(s,\Pi\times\pi)$, equivalently $1/L(s,\pi)$ appears as a degree-$\le n$ polynomial in $q^{-s}$
with generically nonzero coefficients --- a genuine \emph{multi-shell} function. The escape used
here is structural, not a cancellation of $L$-coefficients\textup{:} we never form a multi-shell
function in the first place. The vector $W_0$ is chosen \emph{directly} as the
compactly-supported Kirillov-model vector \eqref{eq:W0}, whose support on $\qtwodiag(G_n,1)$ is the
single bounded set $N_nK_n\subseteq S_1$; this is legitimate precisely because the Kirillov model
$\mathcal K(\Pi,\psi^{-1})$ contains \emph{all} of $\qtwoSw(N_n\backslash G_n;\psi^{-1})$
\textup{(}Lemma~\ref{lem:extend}(b), the Bernstein--Zelevinsky theorem (K1)\textup{)}, so we are
free to prescribe a one-shell function and realize it inside the Whittaker model. The
$\pi$-dependent translate $u_Q$ then acts \emph{only} by the unimodular phase
$\psi(-Q\,g_{nn})$ \textup{(}Lemma~\ref{lem:uniform}(i)\textup{)}, which cannot enlarge the
support. Thus the extra shells are not \emph{cancelled}; they are \emph{absent}, and the collapse
to a monomial is forced by the support of $W_0$ rather than by any arithmetic of the
$L$-factor. This is exactly route \textup{(a)} of the issue: a direct computation of the unfolded
integrand exhibiting single-shell support.

\begin{qtwocorollary}\label{cor:sindep}
With $W_0$ the single $\pi$-independent vector of Lemma~\ref{lem:extend} and $V$ \emph{any}
vector in $\qtwoW(\pi,\psi)$, the problem's Rankin--Selberg integral
\[
\qtwolRS(s,W_0,V):=\int_{N_n\backslash G_n}W_0(\qtwodiag(g,1)u_Q)\,V(g)\,|\det g|^{s-\frac12}\,dg
=I=\int_{K_n}\psi(-Q\,k_{nn})\,V(k)\,dk
\]
is finite and independent of $s$, for every generic $\pi$ and every conductor $c\ge0$. Hence the
Theorem is reduced to producing, for each $\pi$, a single $V\in\qtwoW(\pi,\psi)$ with $I\neq0$.
\end{qtwocorollary}

\begin{proof}
Immediate from Lemma~\ref{lem:unfold}: the integral equals the constant $I$ of \eqref{eq:Idef},
which is an integral of a locally constant function over the compact group $K_n$, hence finite,
and manifestly independent of $s$. Once $I\neq0$ for a suitable $V$, the integral is finite and
nonzero for all $s\in\qtwoC$, with no analytic continuation or pole cancellation required.
\end{proof}

\medskip
\noindent\emph{Remark on the previous $t_Q$ route.} An earlier version of this argument first
right-translated $W_0$ by $t_Q:=u_1d_Q^{-1}=d_Q^{-1}u_Q$ and then attempted to recover the
problem's integral by the substitution $g\mapsto gd_Q^{-1}$, invoking an identity
``$\qtwodiag(g,1)u_Q=\qtwodiag(gd_Q^{-1},1)t_Q$''. That matrix identity is \emph{false}
\textup{(}the two sides differ by a non-unipotent factor\textup{)}, so that route was flawed. The
direct computation above avoids it entirely: we apply the off-mirabolic identity to $u_Q$
\emph{itself} \textup{(}Step~1\textup{)}, which is legitimate because $u_Q\in U$, and never
introduce $d_Q$ or $t_Q$. The outcome is the cleaner constant
$I=\int_{K_n}\psi(-Q\,k_{nn})V(k)\,dk$, supported on the single shell $|\det g|=1$.

\medskip
\noindent\emph{Remark on method.} It is tempting to evaluate $I$ by inserting the
Godement--Jacquet functional equation for the standard $L$-function of $\pi$. That functional
equation, $\gamma(s,\pi,\psi)\,Z(\phi,f,s)=Z(\widehat\phi,f^\vee,1-s)$ with
$f^\vee(g)=f(g^{-1})$, is valid for $f$ a \emph{matrix coefficient} of $\pi$; a Whittaker
function $V$ is not a matrix coefficient, and $g\mapsto V(g^{-1})$ is not a Whittaker function
\textup{(}it carries the wrong $N_n$-equivariance: a left translate of $V$ becomes a right
translate of $g\mapsto V(g^{-1})$\textup{)}, so that functional equation cannot be applied with
$f=V$ and $f^\vee(g)=V(g^{-1})$. We therefore evaluate $I$ directly, using only the defining
properties of the newvector together with the local theory of Whittaker newvectors; the role
classically played by Godement--Jacquet is here played by the \emph{Whittaker-model} computation
of the local $\varepsilon$-factor.

\section*{Nonvanishing of the constant $I$}

By Corollary~\ref{cor:sindep} the Theorem is reduced to the following statement, which we must
establish for \emph{every} generic $\pi$ \textup{(}every conductor $c\ge0$\textup{)}: there is a
$V\in\qtwoW(\pi,\psi)$ with
\[
I(V):=\int_{K_n}\psi(-Q\,k_{nn})\,V(k)\,dk\neq0 .
\]
We give the case $n=1$ in full and reduce $n\ge2$ to the local theory of the essential
\textup{(}new\textup{)} Whittaker vector. In both cases $I(V)$ is a Whittaker-model Gauss-sum
integral whose nonvanishing is the nonvanishing of the local $\varepsilon$-factor.

\subsection*{The case $n=1$ \textup{(}done in full, uniformly in $c$\textup{)}}
Here $\pi$ is a character $\chi$ of $F^\times$ of conductor $\qtwopp^{c}$ \textup{(}$c\ge0$\textup{)},
$G_1=F^\times$, $N_1=\{1\}$, and $\qtwoW(\pi,\psi)=\{a\mapsto \lambda\,\chi(a):\lambda\in\qtwoC\}$.
Take $V(a)=\chi(a)$ \textup{(}so $V(1)=1$\textup{)}, which is $1+\qtwopp^{c}$-invariant because
$\chi$ has conductor $\qtwopp^{c}$. Here $k_{nn}=k$ on $K_1=\qtwooo^\times$, so
\[
I=\int_{\qtwooo^\times}\psi(-Q k)\,\chi(k)\,dk=:G .
\]
\emph{Case $c=0$ \textup{(}$\chi$ unramified, $Q\in\qtwooo^\times$\textup{)}.} Then
$\chi|_{\qtwooo^\times}=1$ and, since $\psi$ is trivial on $\qtwooo$ \textup{(}conductor
$\qtwooo$\textup{)} and $Qk\in\qtwooo$, $\psi(-Qk)=1$; thus
$G=\qtwovol(\qtwooo^\times)=1-\tfrac1q\neq0$.

\emph{Case $c\ge1$.} Then $G$ is a Gauss sum, and we compute $|G|^2$ directly:
\[
|G|^2=\int_{\qtwooo^\times}\!\!\int_{\qtwooo^\times}\chi(k)\overline{\chi(k')}\,
\psi\big(-Q(k-k')\big)\,dk\,dk'
=\int_{\qtwooo^\times}\chi(t)\,J(t)\,dt,\qquad
J(t):=\int_{\qtwooo^\times}\psi\big(-Q(t-1)k'\big)\,dk',
\]
where we substituted $k=k't$ \textup{(}$t\in\qtwooo^\times$\textup{)} and used $|\chi|=1$ on
$\qtwooo^\times$, so $\chi(k)\overline{\chi(k')}=\chi(t)$. For $a\in\qtwoF$,
$\int_{\qtwooo^\times}\psi(ak')\,dk'=\int_{\qtwooo}\psi(ak')\,dk'-\int_{\qtwopp}\psi(ak')\,dk'
=\mathbf 1[\,|a|\le1\,]-\tfrac1q\,\mathbf 1[\,|a|\le q\,]$; with $a=-Q(t-1)$, $|a|=|Q|\,|t-1|=q^{c}|t-1|$,
this gives
\[
J(t)=\mathbf 1[\,t\in1+\qtwopp^{c}\,]-\tfrac1q\,\mathbf 1[\,t\in1+\qtwopp^{c-1}\,].
\]
Therefore
\[
|G|^2=\int_{1+\qtwopp^{c}}\chi(t)\,dt-\tfrac1q\int_{1+\qtwopp^{c-1}}\chi(t)\,dt .
\]
Since $\chi$ is trivial on $1+\qtwopp^{c}$, the first integral is
$\qtwovol(1+\qtwopp^{c})=q^{-c}$. Since $\chi$ has conductor exactly $\qtwopp^{c}$, it is
\emph{nontrivial} on the compact group $1+\qtwopp^{c-1}$, so $\int_{1+\qtwopp^{c-1}}\chi(t)\,dt=0$
\textup{(}integral of a nontrivial character over a compact group\textup{)}. Hence
\[
|G|^2=q^{-c}\neq0,
\]
so $G\neq0$ and $I=G\neq0$, with $|I|=q^{-c/2}=|Q|^{-1/2}$. This settles $n=1$ for every $c\ge0$;
note in particular that the argument is \emph{uniform in $c$}: the conductor enters only through
the modulus-one phase $\psi(-Qk)$ and the conductor characterization of $\chi$, with no upper
bound on $c$.

\subsection*{The case $n\ge2$}
For general $n$ the nonvanishing of $I(V)$ \textup{(}for $V$ the normalized newvector\textup{)} is
the nonvanishing of the local $\varepsilon$-factor of $\pi$, expressed in the Whittaker model. We
isolate the precise input and explain why it holds uniformly in the conductor.

\begin{qtwotheorem}[Whittaker-model $\varepsilon$-factor is nonzero; JPSS81, Matringe
\cite{JPSS81,Matringe}]\label{thm:newvector}
Let $\pi$ be generic irreducible admissible of $G_n$ with conductor $\qtwoqq=\qtwopp^{c}$, let
$V\in\qtwoW(\pi,\psi)$ be the normalized newvector \textup{(}the unique $K_1(\qtwoqq)$-fixed
vector with $V(1)=1$\textup{)}, and let $Q$ generate $\qtwoqq^{-1}$. Then
\[
I=\int_{K_n}\psi(-Q\,k_{nn})\,V(k)\,dk\neq0 .
\]
\end{qtwotheorem}

\noindent\emph{Why this holds, and why it is uniform in $c$.} By Lemma~\ref{lem:unfold},
$I$ is the integral over the maximal compact $K_n$ of the newvector, weighted by the additive
character $k\mapsto\psi(-Q\,k_{nn})$ of conductor exactly $\qtwoqq$ in the $(n,n)$-coordinate
\textup{(}because $|Q|=q^{c}$\textup{)}. This is precisely the Whittaker-model
\textup{(}Gauss-sum\textup{)} avatar of the local $\varepsilon$-factor of $\pi$, and its
nonvanishing is the statement $|\varepsilon(\tfrac12,\pi,\psi)|=1\neq0$. The mechanism is the
local Rankin--Selberg / Godement functional equation \emph{on the Whittaker model}:
\[
Z(1-s,\qtwottilde V,\widehat\Phi)=\gamma(s,\pi,\psi)\,Z(s,V,\Phi),\qquad
Z(s,V,\Phi)=\int_{N_n\backslash G_n}V(g)\,\Phi(e_n g)\,|\det g|^{s}\,dg,
\]
where $\qtwottilde V(g)=V(w_n{}^{t}g^{-1})\in\qtwoW(\qtwottilde\pi,\psi^{-1})$ \textup{(}$w_n$ the
long Weyl element\textup{)}, $\Phi\in\mathcal S(\qtwoF^{n})$, and $\widehat\Phi$ is the
$\psi$-Fourier transform. For the newvector and $\Phi_0=\mathbf 1_{\qtwooo^{n}}$ one has
$Z(s,V,\Phi_0)=L(s,\pi)$; choosing instead the conductor-matching test function
$\Phi_Q$ supported near $Q^{-1}e_n$ converts the right-hand zeta integral into the Gauss sum $I$
above, while the left-hand side is, up to the nonvanishing factor $\gamma(s,\pi,\psi)$ and the
units $L(s,\pi)^{\pm1}$, a nonzero quantity governed by $\varepsilon(\tfrac12,\pi,\psi)$. This is
the content of Matringe's essential-vector computation \cite[Thm.~1.1, \S3]{Matringe}, building on
\cite[\S\S4--5]{JPSS81}; for $n=1$ it specializes to the Gauss-sum computation above. The argument
is uniform in $c$: the newvector exists for every $c$ \textup{(}\cite{JPSS81}\textup{)}, the
functional equation and the $\varepsilon$-factor are defined for every $c$, and
$|\varepsilon(\tfrac12,\pi,\psi)|=1$ holds with no bound on $c$.

\medskip
\noindent We stress that the relevant dual is $\qtwottilde V(g)=V(w_n{}^{t}g^{-1})$, which
\emph{does} lie in the Whittaker model of $\qtwottilde\pi$, in contrast to the spurious
$g\mapsto V(g^{-1})$ that the Godement--Jacquet identity for matrix coefficients would require;
this is why the nonvanishing input is genuinely a Whittaker-model fact and is available uniformly
in the conductor.

\medskip
\noindent\emph{Scope of this subsection.} For the present issue \textup{(}uniformity of a single
$\pi$-independent $W$\textup{)}, the substantive content is Corollary~\ref{cor:sindep}: a
\emph{single} $W_0$, fixed before $\pi$, reduces the problem's integral to the $s$-independent
constant $I(V)$ for every $\pi$ and every conductor $c$, with no boundedness hypothesis. The
remaining nonvanishing $I(V)\neq0$ is proved unconditionally for $n=1$ above; for $n\ge2$ it is
the nonvanishing of the local $\varepsilon$-factor, supplied by the cited essential-vector theory,
and is itself uniform in $c$.

\section*{Conclusion}

\begin{qtwoproposition}\label{prop:final}
Let $W_0\in\qtwoW(\Pi,\psi^{-1})$ be the single $\pi$-independent vector of \eqref{eq:W0} and
Lemma~\ref{lem:extend}. For every generic irreducible admissible $\pi$ of $G_n$, with conductor
$\qtwoqq=\qtwopp^{c}$ \textup{(}any $c\ge0$\textup{)}, generator $Q$ of $\qtwoqq^{-1}$, and
$u_Q=I_{n+1}+QE_{n,n+1}$, there is $V\in\qtwoW(\pi,\psi)$ \textup{(}the normalized newvector\textup{)}
such that, for all $s\in\qtwoC$,
\[
\int_{N_n\backslash G_n} W_0(\qtwodiag(g,1)u_Q)\,V(g)\,|\det g|^{s-\frac12}\,dg
= I=\int_{K_n}\psi(-Q\,k_{nn})\,V(k)\,dk,
\]
a nonzero scalar independent of $s$.
\end{qtwoproposition}

\begin{proof}
By Corollary~\ref{cor:sindep} the integral equals the $s$-independent constant $I$ of
\eqref{eq:Idef} for every $s$ and every $c\ge0$. By the explicit Gauss-sum computation of the case
$n=1$ and Theorem~\ref{thm:newvector} for $n\ge2$, $I\neq0$. Hence the integral is finite and
nonzero for all $s\in\qtwoC$.
\end{proof}

Taking $W:=W_0\in\qtwoW(\Pi,\psi^{-1})$ \textup{(}the single, $\pi$-independent vector of
Lemma~\ref{lem:extend}, constructed before any $\pi$ is named\textup{)} and, for each $\pi$, the
normalized newvector $V\in\qtwoW(\pi,\psi)$, Proposition~\ref{prop:final} shows that the local
Rankin--Selberg integral of the problem is the nonzero constant $I$, hence finite and nonzero for
all $s\in\qtwoC$. Since $W=W_0$ does not depend on $\pi$ \textup{(}nor on its conductor $c$, nor on
$Q$\textup{)}, this is exactly the required uniform statement: one fixed $W$ serves all $\pi$
simultaneously. This proves the Theorem. \qed

\section*{Remarks}

The mechanism is transparent and \emph{functional-equation-free at its core}: the chosen vector
$u_Q W_0$ is supported, after unfolding, on a single $\det$-coset
\textup{(}Lemma~\ref{lem:unfold}\textup{)}, so the Rankin--Selberg integral collapses to a
\emph{constant} $I$ in $s$ rather than to an $L$-factor times a unit. Finiteness and nonvanishing
for all $s$ are therefore equivalent to the single statement $I\neq0$. This nonvanishing is a
local Gauss-sum fact: for $n=1$ it is the classical nonvanishing $|G(\chi,\psi_Q)|=q^{-c/2}$ of a
Gauss sum of a primitive character against an additive character of matching conductor, proved
above in full; for $n\ge2$ it is the corresponding nonvanishing of the newvector
$\varepsilon$-factor integral, supplied by the essential-vector theory of \cite{JPSS81,Matringe}.

We emphasize the methodological point underlying the present write-up. The natural temptation is
to evaluate $I$ via the Godement--Jacquet functional equation for the standard $L$-function of
$\pi$. That functional equation is a statement about \emph{matrix coefficients} of $\pi$, and its
dual is $f^\vee(g)=f(g^{-1})$; for a Whittaker function $V$ the function $g\mapsto V(g^{-1})$ is
\emph{not} a Whittaker function, so the Godement--Jacquet identity does not apply to $f=V$. The
correct functional equation in the Whittaker model is the Jacquet--Piatetski-Shapiro--Shalika one,
whose dual $\qtwottilde V(g)=V(w_n{}^{t}g^{-1})$ \emph{is} a Whittaker function \textup{(}of
$\qtwottilde\pi$\textup{)}. Accordingly, all functional-equation input here is on the Whittaker
side, and the load-bearing nonvanishing is reduced to the explicit local computation of $I$.

The $n=1$ case recovers the normalized Gauss sum that appears in subconvexity and moment
applications \cite{Sarnak,MV}; the general-$n$ statement, with a single unipotent translate rather
than an average, is the present construction. A version with an average over many translates
appears in \cite{BKL}.

\begin{thebibliography}{9}
\bibitem{Bernstein} J.~N.~Bernstein.
  \emph{$P$-invariant distributions on $\qtwoGL(N)$ and the classification of
  unitary representations of $\qtwoGL(N)$ (non-archimedean case)}. In
  \emph{Lie group representations II}, LNM \textbf{1041}, Springer, 1984,
  50--102.
\bibitem{BKL} A.~R.~Booker, M.~Krishnamurthy, M.~Lee.
  \emph{Test vectors for Rankin--Selberg $L$-functions}.
  J.\ Number Theory \textbf{209} (2020), 37--48.
\bibitem{GJ72} R.~Godement, H.~Jacquet.
  \emph{Zeta functions of simple algebras}. LNM \textbf{260}, Springer, 1972.
\bibitem{JPSS81} H.~Jacquet, I.~I.~Piatetski-Shapiro, J.~Shalika.
  \emph{Conducteur des repr\'esentations du groupe lin\'eaire}.
  Math.\ Ann.\ \textbf{256} (1981), 199--214.
\bibitem{JPSS83} H.~Jacquet, I.~I.~Piatetski-Shapiro, J.~Shalika.
  \emph{Rankin--Selberg convolutions}.
  Amer.\ J.\ Math.\ \textbf{105} (1983), 367--464.
\bibitem{Matringe} N.~Matringe.
  \emph{Essential Whittaker functions for $\qtwoGL(n)$}.
  Doc.\ Math.\ \textbf{18} (2013), 1191--1214.
\bibitem{MV} P.~Michel, A.~Venkatesh.
  \emph{The subconvexity problem for $\qtwoGL_2$}.
  Publ.\ Math.\ IH\'ES \textbf{111} (2010), 171--271.
\bibitem{Sarnak} P.~Sarnak.
  \emph{Fourth moments of Gr\"ossencharakteren zeta functions}.
  Comm.\ Pure Appl.\ Math.\ \textbf{38} (1985), 167--178.
\end{thebibliography}

\endgroup

\endgroup
